\section{Ablation Study}
\label{sec:ablation}

The main results (Tables~\ref{tab:perproblem}--\ref{tab:llama-sweep})
use the default \emph{strategy-enumerate} Phase~3 loop. The
ablations below show (i)~that two alternative LLM-driven Phase~3
styles (cc-react, multi-island GA) reach the same deployed
end-to-end performance but cost $\sim$25\% more wall-time and
$\sim$20\% more tokens (Table~\ref{tab:ablation-cost}), and
(ii)~that disabling the simulator's per-op cost feedback collapses
the $1.40\times$ OLMoE headline to $1.00\times$ because the
LLM-judge falls back to whichever pattern wins a small-shape
microbench --- structurally the inline developer baseline.

\subsection{Alternative Phase-3 loops: cc-react and multi-island GA}
\label{sec:ablation-styles}

\textbf{cc-react} is a single-trajectory ReAct agent on Claude
Code~\cite{anthropic_claude_code} that maintains a scratchpad of
running champion plus per-op cost feedback and proposes
sequential refinements. \textbf{multi-island GA} uses three or
more populations seeded from distinct baselines with LLM-mutation
per round and occasional champion exchange, modelled on
AlphaEvolve~\cite{novikov2025alphaevolve} with our simulator as
fitness. Both share Phase~1's simulator and Phase-4/5's gates.

\textbf{Result}: per-problem cells sit within $\sim$10\% of
strategy-enumerate at every scope
(Appendix~Table~\ref{tab:ablation-perproblem}); OLMoE end-to-end
reproduces $\mathbf{1.40\times}$ within $<0.5\%$
(Appendix~Table~\ref{tab:ablation-e2e}); Llama amp1--amp4 stays
within $\sim$10\% (Appendix~Table~\ref{tab:ablation-llama-sweep}). All three
styles converge to the same deployed Phase-5 winner. The Phase~3
search shape does not move the headline.

\textbf{Why strategy-enumerate is preferred}: it gets to the same
deployed strategy with $\sim$25\% less wall-time and $\sim$20\%
less LLM token spend (Table~\ref{tab:ablation-cost}) at the same
total LLM-call budget of $10$ calls per problem, distributed as
follows: one call asks the LLM to \emph{enumerate} $K{=}5$
structurally distinct strategy sketches, $K$ calls implement each
sketch, and the remaining $2R{=}4$ calls (two rounds, two
candidates) refine the top two implementations against the
simulator. The first six calls thus spend the budget on
\emph{exploring different ideas}, and only the last four refine
within the best ones. The classical greedy-vs-portfolio
trade-off~\cite{kocsis2006ucb, russell2010ai} applies: an LLM's
first emission for one template explores a narrower neighbourhood
than one round each on $K$ distinct templates, which is why
cc-react and multi-island spend more rounds to converge. Prompt
scaffolds for the alternatives are in
Appendix~\ref{sec:search-prompts}; per-style OLMoE end-to-end
and Llama amp1--amp4 cells are in
Appendix~Tables~\ref{tab:ablation-e2e},~\ref{tab:ablation-llama-sweep}.

\begin{table}[!htbp]
\centering\small
\setlength{\tabcolsep}{4pt}
\begin{tabular}{lrr}
\toprule
\textbf{Phase-3 style} & \textbf{Wall} & \textbf{LLM cost} \\
                       & (min) & (Sonnet 4.5) \\
\midrule
strategy-enumerate (default) & 56 & \textbf{\$19} \\
cc-react                     & 75 & \$22 \\
multi-island                 & 78 & \$24 \\
\bottomrule
\end{tabular}
\caption{Per-style search cost over the eight collective problems
on the 7-node cluster. Strategy-enumerate is the cheapest of the
three by both wall-time and LLM token cost; cc-react and
multi-island spend extra rounds without changing the deployed
strategy's end-to-end performance
(Appendix~Tables~\ref{tab:ablation-e2e},~\ref{tab:ablation-llama-sweep}).}
\label{tab:ablation-cost}
\end{table}

\subsection{No-simulator ablation: necessity of Phase-3 cost feedback}
\label{sec:ablation-nosim}

To isolate the simulator's per-op cost feedback, strategy-enumerate
runs in \texttt{--no-simulator} mode on the same eight-problem set
under the same per-style LLM budget. Phase~1 still builds the
simulator, but its output is hidden from the agent in Phase~3;
the saved Phase-3-refinement budget is spent letting the agent
enumerate candidates, write its own microbenchmarks, and produce
code. Phase~5 cannot rank by simulator either, so all
Phase-4-shape-gate survivors are shown back to the LLM with their
HW microbench latencies and code, and the LLM emits the candidate
name it considers best (``LLM-judge'' Phase~5). Search wall-time
is unchanged ($\sim$59 vs $\sim$56 min across the eight problems).

\paragraph{What got deployed and end-to-end outcome.}
With the simulator hidden, the LLM-judge ranks survivors by h7-style
HW microbench latencies and falls back to whichever pattern wins at
small shapes. On all three load-bearing OLMoE primitives this is
structurally the inline developer baseline: for \texttt{alltoallv}
and \texttt{uniform\_a2a} the judge picks the
\texttt{allgather\_reduce\_scatter} (AG+T+RS) pattern; for
\texttt{dxe} it picks the \texttt{full\_vocab\_ag} 1-AG strategy,
both call-level bit-equivalent to the inline baseline the OLMoE
harness exercises. The end-to-end OLMoE ratio is therefore
$\mathbf{1.00\times}$ (Table~\ref{tab:ablation-nosim}) ---
exactly what the mechanistic reading predicts when both paths run
bit-identical collective strategies for the load-bearing
primitives, and the $1.40\times$ headline collapses entirely. The
per-problem deployments and a per-problem cell table are in
Appendix~\ref{sec:perproblem-ablation-tables} (Table~\ref{tab:ablation-nosim-perproblem}).

\begin{table}[!htbp]
\centering\footnotesize
\setlength{\tabcolsep}{3pt}
\begin{tabular}{lrrr}
\toprule
\textbf{Harness} & \textbf{base (ms)} & \textbf{no-sim (ms)} & \textbf{Ratio} \\
\midrule
OLMoE-10B (a2av + dxe)     & 4399.8 & 4393.4 & $\mathbf{1.00\times}$ \\
Llama-7B$^{\diamondsuit}$  & 76.9   & 76.0   & $1.01\times$ \\
\bottomrule
\end{tabular}
\caption{No-simulator ablation end-to-end. OLMoE $\mathbf{1.00\times}$
is the central finding: stripping simulator scores collapses the
$1.40\times$ headline. $^{\diamondsuit}$Llama-7B drops PP, uses
TP$+$FSDP$+$lbar; per-problem rows in
Appendix~\ref{sec:perproblem-ablation-tables}.}
\label{tab:ablation-nosim}
\end{table}

\subsection{Cluster-size generalization: 7-node to 3-node}
\label{sec:ablation-3node}

The deployed strategies have a few \emph{topology constants}
baked in at codegen --- specifically, the total number of ranks,
the half-cluster size used by the pipeline-parallel split, and
the number of physical nodes. On the 7-node cluster these
constants are $(\text{total ranks}, \text{half}, \text{nodes})
{=}(224, 112, 7)$. To check that the agent's strategies are not
overfitted to this specific cluster shape --- a real risk for
any strategy calibrated against device-specific characteristics
--- we redeploy the same strategies on smaller clusters by
setting these constants to $(96, 48, 3)$ at 3 nodes and
$(160, 80, 5)$ at 5 nodes, and re-derive the model-shape
parameters (e.g.\ MoE expert count and Llama hidden dimension)
so the model still fits cleanly at the new cluster size
(Appendix~\ref{sec:3node-detail}). If the strategies were
overfitted to the 7-node shape, the speedup would either
disappear or invert at these smaller clusters. Table~\ref{tab:ablation-3node}
shows the speedup persists in both directions, indicating that
the agent's strategies capture device behavior that
generalizes across cluster sizes rather than tuning to one
particular topology.

\begin{table}[!htbp]
\centering\footnotesize
\setlength{\tabcolsep}{4pt}
\begin{tabular}{lrrr}
\toprule
\textbf{Harness} & \textbf{base (ms)} & \textbf{agent (ms)} & \textbf{Speedup} \\
\midrule
OLMoE-10B 7n (224r) & 4404.3 & 3139.8 & $1.40\times$ \\
OLMoE-8B 5n (160r)  & 3258.0 & 2173.0 & $\mathbf{1.50\times}$ \\
OLMoE-5B 3n (96r)   & 2547.4 & 1422.0 & $\mathbf{1.79\times}$ \\
\midrule
Llama-7B 7n (224r)  & 76.90 & 21.50 & $3.58\times$ \\
Llama-7B 5n (160r)  & 67.94 & 24.77 & $\mathbf{2.74\times}$ \\
Llama-7B 3n (96r)   & 67.82 & 27.20 & $\mathbf{2.49\times}$ \\
\bottomrule
\end{tabular}
\caption{Cluster-size generalization: same runtimes re-deployed at
5-node and 3-node with topology constants re-fitted. OLMoE speedup
widens as cluster shrinks; Llama speedup compresses. See
§\ref{sec:ablation-3node}.}
\label{tab:ablation-3node}
\end{table}

